\documentclass[10pt,letterpaper]{article}
\usepackage[margin=0.8in]{geometry}
\usepackage{amsmath,amssymb,amsthm,mathtools}
\usepackage[T1]{fontenc}
\usepackage{lmodern,microtype,booktabs,xcolor,enumitem,fancyhdr}
\usepackage[colorlinks=true,linkcolor=blue!45!black,urlcolor=blue!45!black]{hyperref}
\usepackage{listings}
\definecolor{ink}{HTML}{17324D}
\definecolor{muted}{HTML}{566474}
\pagestyle{fancy}\fancyhf{}
\fancyhead[L]{\small\color{muted}Mixed-integer Gr\"unbaum bound}
\fancyhead[R]{\small\color{muted}7 November 2026}
\fancyfoot[C]{\small\thepage}
\setlength{\headheight}{14pt}
\setlength{\parindent}{0pt}\setlength{\parskip}{5pt}
\setlist[itemize]{leftmargin=1.4em,itemsep=3pt,topsep=3pt}
\newtheorem{theorem}{Theorem}
\newtheorem{lemma}{Lemma}
\newcommand{\R}{\mathbb R}\newcommand{\Z}{\mathbb Z}
\newcommand{\conv}{\operatorname{conv}}\newcommand{\vol}{\operatorname{vol}}
\newcommand{\MI}{\mu_{\mathrm{MI}}}\newcommand{\depth}{\operatorname{depth}}
\newcommand{\eps}{\varepsilon}\newcommand{\Ga}{\Gamma}
\newcommand{\head}[1]{\vspace{5pt}{\large\bfseries\color{ink}#1}\par}
\lstset{basicstyle=\ttfamily\footnotesize,columns=fullflexible,keepspaces=true,
  frame=single,rulecolor=\color{black!20},xleftmargin=5pt,xrightmargin=5pt,
  aboveskip=7pt,belowskip=5pt,showstringspaces=false}
\hypersetup{pdftitle={Counterexample to the 1/(e 2^n) mixed-integer Grunbaum bound},
 pdfauthor={Hongyu Cheng and Phillip Kerger}}
\begin{document}
\thispagestyle{plain}
{\LARGE\bfseries\color{ink}A counterexample to the $1/(e2^n)$\\[3pt]
mixed-integer Gr\"unbaum bound\par}
\vspace{5pt}
{\color{muted}Hongyu Cheng and Phillip Kerger\hfill 7 November 2026}

For a bounded Borel set $D\subset\R^n\times\R^d$, write
\[
 \MI(D)=\sum_{z\in\Z^n}\vol_d\{x:(z,x)\in D\},\qquad
 \depth_C(y)=\inf_{\substack{H\text{ closed halfspace}\\y\in H}}
                    \frac{\MI(C\cap H)}{\MI(C)}.
\]
A halfspace is proper: its defining linear functional is nonzero. All volumes
below are ordinary Lebesgue volumes in the indicated continuous coordinates.

\begin{theorem}\label{thm:main}
There is a full-dimensional rational polytope
$C\subset\R^{76}\times\R^d$, where $d=1{,}372{,}131+2^{93}$, such that
$S=C\cap(\Z^{76}\times\R^d)$ is nonempty, every nonempty integer fiber has
nonempty $d$-dimensional interior, and
\[
                 \sup_{y\in S}\depth_C(y)<\frac{1}{e2^{76}}.
\]
Consequently the specific $1/e$ bound and the larger sharp
$\Ga_d=(d/(d+1))^d$ bound both fail.
\end{theorem}
This does not disprove a bound $c\,2^{-n}$ for some smaller positive absolute
constant $c$. The construction and strict mixed-integer-volume conclusion
have been verified in Lean; the precise scope is recorded on page~\pageref{sec:formal}.

\head{The idea of the construction}
There will be just one exceptional integer fiber, over the origin. Every
other integer base will be exposed: a linear functional on the integer
coordinates uniquely maximizes there. A sufficiently large multiple of that
functional allows us to isolate a continuous cap in its fiber. We therefore
have two separate tasks: make every point over the origin shallow, and
provide suitable caps at all the exposed bases.

Before a final simplex factor is added, the intended fibers have these volumes:
\begin{center}
\begin{tabular}{@{}lrr@{}}\toprule
Integer bases & Number & Volume of each fiber\\\midrule
Origin & $1$ & $b$\\
Nonzero vertices of the unit cube & $2^{76}-1$ & $bm$\\
Negative coordinate units & $76$ & $b\alpha$\\\bottomrule
\end{tabular}
\end{center}
The cube matches the factor $2^{76}$ in the proposed bound. The additional
negative-unit fibers create a surplus in the total mass:
\[
 \MI(C_0)=bW,\qquad W=2^{76}m+\delta,\qquad\delta=1+76\alpha-m.
\]
We shall construct four halfspaces that cover the entire origin fiber. Each
contains a fraction $27/64$ of that fiber, exactly $19$ outer fibers, and
no other fibers. Its mass is consequently $bL$, where
$L=27/64+19\alpha$. The quantitative goals are
\[
                 0<\alpha\le m,\qquad\delta>0,\qquad L<m/e.
\]
The first two make each exposed fiber smaller than $2^{-76}$ of the total;
the third deals with the origin. A common $\ell$-simplex factor then supplies
caps of fraction $\Ga_\ell\to1/e$ at the exposed bases, without changing the
origin-cut ratios. For sufficiently large finite $\ell$, the uniform bound is
\[
              \frac{\max\{L,m\Ga_\ell\}}{2^{76}m+\delta}
                       <\frac1{e2^{76}}.
\]
The proof below realizes this scheme. The main compatibility issue is that
taking a convex hull must not enlarge the origin fiber. A grouped Minkowski
containment will resolve exactly this issue. We keep the parameters symbolic
until their geometric roles are clear, and give exact rational values in
Section~4.

\newpage
\head{1. Four source types and the two continuous factors}
Let $B_0\subset\R^3$ be the standard tetrahedron translated to have centroid
zero, with barycentric coordinates $\lambda_0,\ldots,\lambda_3$.
Thus $f_j=\lambda_j-1/4$ is linear and $\sum_j f_j=0$.
For a small rational $\eps>0$ and $a\in\{0,1,2,3\}$ define
\[
 T_a=\{x\in B_0:\lambda_j(x)\le1/4\ (j\ne a)\},\qquad
 \lambda_j(c_a)=\begin{cases}1/4+3\eps&j=a,\\1/4-\eps&j\ne a.\end{cases}
\]
The three coordinates with $j\ne a$ identify $T_a$ with $[0,1/4]^3$ and
$B_0$ with the standard simplex in those coordinates. Hence
\[
 \theta:=\frac{|T_a|}{|B_0|}=\frac{(1/4)^3}{1/6}=\frac3{32},\qquad
 \frac{|B_0\cap\{f_j\ge0\}|}{|B_0|}=\left(\frac34\right)^3=\frac{27}{64}.
 \tag{1}
\]
The second identity follows by homothety about vertex $j$. The reason for
these sources is that $T_a$ lies in $\{f_j\le0\}$ for all $j\ne a$:
a source of type $a$ can be excluded by three of the four cuts. Its anchor
$c_a$ lies just inside these three planes. Contraction about this anchor
will give strict separation.

Choose rational parameters and an integer $h\ge1$ satisfying
\[
\begin{gathered}
 0<\eps<1/4,\quad 1/2<s<1,\quad 1/2<q<1,\quad u=2(1-q),\\
 0<t<\eta<2(1-s),\qquad
 t(1/4+3\eps)\le2(1-s)(1/4-\eps).
\end{gathered}\tag{2}
\]
The box $F=[-1,1]^h$ will tune volumes; its coordinates never enter the four
origin cuts. Define
\[
 \begin{aligned}
 B&=B_0\times F,& P_a&=[c_a+s(T_a-c_a)]\times qF,\\
 A_a&=t(T_a-c_a)\times uF,& K(a,k)&=2P_a+(k-1)(c_a,0)\quad(k\ge1).
 \end{aligned}\tag{3}
\]
All these sets are rational polytopes with nonempty interior; $c_a$ is an
interior point of $T_a$. Writing $b=|B|>0$, scaling and product volumes give
\[
 |K(a,k)|=bm,\quad |A_a|=b\alpha,\qquad
 m=(2s)^3\theta(2q)^h,\quad\alpha=t^3\theta u^h.\tag{4}
\]

Since $2q>1$ and $u<1$, increasing $h$ increases $m$ and decreases $\alpha$.
The relation $u=2(1-q)$ keeps these changes compatible with convexification.
The later simplex factor has a different job: it supplies small containing
caps at exposed bases, multiplying every existing volume by the same factor.

\head{2. The compatibility needed at the origin}
\begin{lemma}[Grouped containment]\label{lem:group}
For $k\ge1$ and arbitrary source types $a,a_1,\ldots,a_k$,
\[
                  K(a,k)+\sum_{r=1}^k A_{a_r}\subset (k+1)B.\tag{5}
\]
\end{lemma}
\begin{proof}
Set $c_-=1/4-\eps$ and $c_+=1/4+3\eps$. The final assumption in (2) is
\[
                         tc_+\le2(1-s)c_-.\tag{6}
\]
For $p\in c_a+s(T_a-c_a)$ and $v_r\in t(T_{a_r}-c_{a_r})$ we have
$\lambda_j(p)\ge(1-s)\lambda_j(c_a)$ and
$f_j(v_r)\ge-t\lambda_j(c_{a_r})$. Therefore, for
$y=(2p+(k-1)c_a+\sum_r v_r)/(k+1)$,
\[
 \begin{aligned}
 (k+1)\lambda_j(y)
 &\ge (k+1-2s)c_- -ktc_+\\
 &\ge \bigl[k+1-2s-2k(1-s)\bigr]c_-
   =(k-1)(2s-1)c_-\ge0.
 \end{aligned}
\]
Thus $y\in B_0$. The free-coordinate radius is at most $2q+ku\le k+1$,
since $k+1-2q-ku=(k-1)(2q-1)\ge0$. These two statements prove (5).
\end{proof}

\newpage
\head{The convex body and its actual fibers}
Index the integer coordinates by
$I=\{0,1,2,3\}\times\{1,\ldots,19\}$; let $\tau(i)$ be the first coordinate.
For each nonempty $E\subset I$, fix once and for all one member,
let $a(E)$ be that member's type, and set $z_E=\sum_{i\in E}e_i$.
The choice $a(E)\in\tau(E)$ ensures that if $E$ has no coordinate of type
$j$, then $a(E)\ne j$; cut $j$ will therefore exclude its source. The translation $(|E|-1)c_{a(E)}$ in (3) makes
the corresponding negative continuous score grow with $|E|$.
Define $C_0\subset\R^{76}\times\R^{3+h}$ to be the convex hull of
\[
 \{0\}\times B,\qquad
 \{z_E\}\times K(a(E),|E|)\quad(\varnothing\ne E\subset I),\qquad
 \{-e_i\}\times A_{\tau(i)}\quad(i\in I).\tag{7}
\]
Why is (5) the relevant condition? The projected generators satisfy
\[
                  0=\frac{z_E+\sum_{i\in E}(-e_i)}{|E|+1}.
\]
Such groups can contribute to the origin fiber when we take the convex hull.
Lemma~\ref{lem:group} says their continuous averages remain inside $B$.
The next argument shows that every contribution to the origin can be
regrouped in precisely this way.

\begin{lemma}[Exact integer fibers]\label{lem:fibers}
The only integer bases of $C_0$ are $0$, the nonzero cube vertices $z_E$,
and the negative units $-e_i$. Their fibers are exactly the sets in (7).
The polytope $C_0$ has nonempty ambient interior.
\end{lemma}
\begin{proof}
Every coordinate of its projection lies in $[-1,1]$. If an integer projected
point has a negative coordinate, that coordinate is $-1$, which forces every
positive-weight generator in a convex representation to be the unique
negative unit with that coordinate. Otherwise the point is a cube vertex.
For nonempty $E$, the functional with coefficients $1$ on $E$ and $-1/2$
outside $E$ uniquely exposes $z_E$ among all bases in (7): its value there
is $|E|\ge1$, whereas negative units have value at most $1/2$.
The functional $-z_i$ uniquely exposes $-e_i$. Their fibers are therefore exact.

At the origin, aggregate a convex representation into weights $w_0,w_E,v_i$
and points $x_0\in B$, $x_E\in K(a(E),|E|)$, $y_i\in A_{\tau(i)}$.
Convexity of the assigned fibers permits this aggregation. Zero integer
coordinate and total weight one imply
\[
 v_i=\sum_{E\ni i}w_E,\qquad w_0+\sum_E(|E|+1)w_E=1.
\]
The continuous point becomes
\[
 w_0x_0+\sum_Ew_E\left(x_E+\sum_{i\in E}y_i\right)\in B
\]
by Lemma~\ref{lem:group}; unused zero-weight points may be chosen arbitrarily.
The reverse inclusion follows from (7). Finally the origin fiber spans all
continuous directions, and the projected bases span all integer directions.
Thus the affine hull of $C_0$ is the whole ambient space. A full-dimensional
convex polytope has nonempty interior.
\end{proof}

\head{3. Four small cuts cover the origin fiber}
Let $R=231$. For each $j\in\{0,1,2,3\}$ define the proper closed halfspace
\[
 H_j=\left\{(z,x,w): f_j(x)+\sum_{i\in I}\beta_i z_i\ge0\right\},\qquad
 \beta_i=\begin{cases}-R&\tau(i)=j,\\\eta\eps&\tau(i)\ne j.\end{cases}\tag{8}
\]
Here $x\in\R^3$ and $w\in\R^h$. Call a coordinate with $\tau(i)=j$
\emph{paid}: its outer fiber will be included by this cut. The continuous functional $f_j$ is nonzero. Our assumptions give
\[
             0<t<\eta<2(1-s)<1,\qquad \eta\eps<1.\tag{9}
\]
\newpage
For $j\ne a$, the definitions give
\[
 \begin{split}
 f_j(P_a)&\le-(1-s)\eps,\qquad f_j(A_a)\le t\eps,\\
 f_j(K(a,k))&\le-(k+1-2s)\eps.
 \end{split}\tag{10}
\]
These inequalities mean pointwise upper bounds on the indicated sets,
ignoring their free coordinates. For every $a,j$, also
$|f_j(A_a)|\le2t$ and $f_j(K(a,k))\le k+1$: the latter follows from
$P_a\subset B_0\times\R^h$, $c_a\in B_0$, and $|f_j|\le1$ on $B_0$.
There are two ways to exclude a cube fiber: a paid coordinate contributes
a large negative integer term, or its source contributes a negative
continuous term. Using the exact census, the complete accounting is:
\begin{itemize}
\item At an unpaid negative unit, the score in (8) is at most
$(t-\eta)\eps<0$. At a paid negative unit it is at least $R-2t>0$,
so precisely the $19$ paid outer fibers are included completely.
\item If $E$ contains a paid coordinate and $k=|E|$, its entire fiber has
score at most $(k+1)-R+(k-1)\eta\eps\le2k-R<0$, since $k\le76$.
\item If $E$ contains no paid coordinate, then $a(E)\ne j$, so its score
is at most $[k\eta-(k+1-2s)]\eps$. By (9) this is strictly less than
$-(k-1)(2s-1)\eps\le0$, including when $k=1$.
\end{itemize}
Thus every nonzero cube fiber is strictly excluded. At the origin, (8)
reduces to $f_j\ge0$. Consequently (1) and (4) give
\[
             \MI(C_0\cap H_j)=bL,\qquad L:=\frac{27}{64}+19\alpha.\tag{11}
\]
Since $\sum_jf_j=0$, every point of the origin fiber belongs to at least
one of these four cuts. All halfspaces are closed; no boundary mass has
been removed in this calculation.

\head{4. Choosing the parameters to balance the masses}
We now specify the rational values promised in (2):
\[
\begin{gathered}
 h=1{,}372{,}128,\quad \eps=10^{-10},\quad
 s=\frac{5000001}{10^7},\quad t=\frac{9999997}{10^7},\\
 \eta=\frac{3999999}{4000000},\quad q=\frac{500001}{10^6},\quad u=2(1-q).
\end{gathered}
\]
The elementary inequalities in (2) hold. For the less immediate one,
\[
                  2(1-s)c_- -tc_+=\frac{2460000011}{10^{17}}>0.
\]
Recall the mass surplus and total, now justified by the exact fiber census:
\[
 \delta=1+76\alpha-m,\qquad W=2^{76}m+\delta,\qquad \MI(C_0)=bW.\tag{12}
\]
The exact power certificates give the following strict rational intervals:
\[
 \begin{aligned}
 m_-:=\frac{1458094201550}{10^{12}}&<m<
                 \frac{1458094201551}{10^{12}}=:m_+,\\
 \alpha_-:=\frac{6027739233}{10^{12}}&<\alpha<
                 \frac{6027739234}{10^{12}}=:\alpha_+.
 \end{aligned}\tag{13}
\]
Page~\pageref{sec:certificate} supplies a short exact-arithmetic reproduction
and its enclosure invariant. In particular $0<\alpha<m$.
Put $U=27182819/10^7$. The exponential series gives
\[
 e<\sum_{j=0}^{10}\frac1{j!}+\frac{12}{11\cdot11!}<U;
\]
the tail starting with $1/11!$ is bounded by a geometric series of ratio $1/12$.

The margins required below are simple rational subtractions:
\[
 \begin{aligned}
 \frac{m_-}{U}-\frac{27}{64}-19\alpha_+
     &=\frac{11151455803863}{13591409500000000000}>0,\\
 1+76\alpha_- -m_+-\frac{m_+}{2^{17}}
     &=\frac{374312936753}{131072000000000000}>0.
 \end{aligned}\tag{14}
\]
\newpage
Combining (13)--(14) with $e<U$ proves
\[
                       L<\frac me,\qquad \delta>\frac m{2^{17}}>0.\tag{15}
\]

\head{5. Why a finite simplex factor completes the example}
The surplus $\delta$ has a specific purpose: it lets us use a simplex cap
slightly larger than $1/e$. The required comparison is
$e\Ga_\ell<1+\delta/(2^{76}m)$. The estimate $\delta/m>2^{-17}$ suggests
$\ell=2^{76+17}$; this explains the large exponent rather than treating it
as an unexplained parameter.
Let $J\subset\R^\ell$ be the standard simplex, with $\ell=2^{93}$, and
set $C=C_0\times J$, grouping all continuous coordinates together. Its actual
integer fibers are precisely the fibers of Lemma~\ref{lem:fibers} times $J$.
It is a full-dimensional rational polytope and all its nonempty fibers have
full dimension. Its mass is $bW|J|>0$.

\begin{lemma}[Isolating a simplex cap]\label{lem:padding}
At every mixed-integer point of $C$ over a nonzero base, a proper closed ambient
halfspace contains the point and has mass at most $bm\Ga_\ell|J|$, where
$\Ga_\ell=(\ell/(\ell+1))^\ell$.
\end{lemma}
\begin{proof}
Let $z_*$ be its base and $v\in J$ its simplex coordinate. Some barycentric
coordinate $\rho_j(v)$ is at least $a=1/(\ell+1)$. The cap
$\{v'\in J:\rho_j(v')\ge a\}$ is the homothetic image
$a v_j+(1-a)J$, where $v_j$ is the corresponding vertex. Its volume is
$\Ga_\ell|J|$.

By Lemma~\ref{lem:fibers}, a linear functional $g$ on integer coordinates
strictly exposes $z_*$. There are only finitely many other bases. Thus
$\Delta=\min_{z\ne z_*}[g(z_*)-g(z)]>0$. Choose $M>1/\Delta$. The halfspace
\[
        \rho_j(v')-a+M[g(z')-g(z_*)]\ge0\tag{16}
\]
is proper, contains the given point, cuts out exactly that cap at $z_*$, and
misses every other fiber, since $\rho_j(v')\le1$ on $J$.
Its affine expression is equivalent to a linear functional with a real
threshold; the continuous linear part of $\rho_j$ is nonzero.
The old fiber at $z_*$ has volume either $bm$ or $b\alpha\le bm$.
Multiplication of product volumes proves the claim.
\end{proof}

For every positive integer $k$,
\[
                  e\left(\frac{k}{k+1}\right)^k<1+\frac1k.\tag{17}
\]
Indeed $\log(1+1/k)>1/(k+1)$, obtained by integrating $1/x$ on
$[1,1+1/k]$. Thus $(k+1)\log(1+1/k)>1$, which is equivalent to (17).
By (15) and $\ell=2^{76+17}$,
\[
       e\Ga_\ell<1+\frac1\ell
       <1+\frac{\delta}{2^{76}m}=\frac{W}{2^{76}m}.\tag{18}
\]

\begin{proof}[Proof of Theorem~\ref{thm:main}]
At every origin-fiber point, extend a containing cut (8) without using the
$J$ coordinates; (11) gives normalized mass $L/W$. All other points are
covered by Lemma~\ref{lem:padding}, with normalized mass at most
$m\Ga_\ell/W$. By (15), $W>2^{76}m$, and (18), their uniform upper bound is
\[
       r:=\frac{\max\{L,m\Ga_\ell\}}{W}<\frac1{e2^{76}}.\tag{19}
\]
Therefore $\sup_{y\in S}\depth_C(y)\le r$. There is no limiting construction
or assumption that a deepest point exists. The origin belongs to $S$.
All fibers have full dimension, so $\dim_H S=d$. Euclidean $\mathcal H_d$,
normalized to agree with Lebesgue measure on $d$-dimensional affine planes,
is exactly the mixed-integer measure on this finite union of disjoint fibers. Finally
$\Ga_d>1/e$, since $d\log(1+1/d)<1$. This also refutes the sharp bound.
\end{proof}

\newpage
\head{Appendix A. Reproducing the finite arithmetic}\label{sec:certificate}
The large powers in (4) require no floating-point computation. Set $Q=10^{30}$.
For nonnegative rational intervals, multiply their endpoints and round the
lower endpoint down and the upper endpoint up to multiples of $1/Q$.
Positivity preserves the enclosure at every multiplication; induction through
binary exponentiation therefore gives a rigorous enclosure of each power.

The following standalone Python code (standard library only) verifies (13)
and the two margins. Every assertion is an exact rational comparison. The
Lean certificate independently proves its individual rational inequalities
inside the kernel; this Python program is not part of Lean's trusted base.
\begin{lstlisting}[language=Python]
from fractions import Fraction as F
from math import floor, ceil, factorial
Q = 10**30

def mul(A, B):
    return (F(floor(Q*A[0]*B[0]), Q),
            F(ceil(Q*A[1]*B[1]), Q))

def power(x, n):
    out, base = (F(1), F(1)), (x, x)
    while n:
        if n % 2: out = mul(out, base)
        n //= 2
        if n: base = mul(base, base)
    return out

s, t, q = F(5000001,10**7), F(9999997,10**7), F(500001,10**6)
h, theta = 1372128, F(3,32)
G, D = power(2*q,h), power(2*(1-q),h)
M = tuple((2*s)**3*theta*x for x in G)
A = tuple(t**3*theta*x for x in D)
ml, mh = F(1458094201550,10**12), F(1458094201551,10**12)
al, ah = F(6027739233,10**12), F(6027739234,10**12)
assert ml < M[0] <= M[1] < mh
assert al < A[0] <= A[1] < ah
U = F(27182819,10**7)
assert sum((F(1,factorial(j)) for j in range(11)), F(0)) \
       + F(12,11*factorial(11)) < U
assert ml/U - F(27,64) - 19*ah == \
       F(11151455803863,13591409500000000000)
assert 1 + 76*al - mh - mh/2**17 == \
       F(374312936753,131072000000000000)
\end{lstlisting}

\head{Appendix B. Formal verification and scope}\label{sec:formal}
The final unconditional Lean theorem is
\begin{center}\small\ttfamily
SharpOC.BoxCounterexample.one\_over\_e\_counterexample
\end{center}
in \texttt{lean/SharpOC/BoxMeasure.lean}. It proves the stated dimensions,
an actual compact convex body with nonempty ambient interior, positive finite
mass summed over \emph{all} integer bases, full-dimensional nonempty fibers,
and one uniform strict bound below $1/(e2^{76})$ using proper closed ambient
halfspaces at every actual mixed-integer point.

The pinned project uses Lean 4.34.1 and mathlib commit
\href{https://github.com/leanprover-community/mathlib4/commit/d13f23b723b8a846827a245b89c10fc7d3f11612}{\texttt{d13f23b723b8}}. Both
\texttt{lake build} and \texttt{lake env lean CheckAxioms.lean} passed.
The final theorem depends only on \texttt{propext}, \texttt{Classical.choice},
and \texttt{Quot.sound}; there are no admitted proofs or extra axioms.
The formal theorem uses canonical Lebesgue measure in continuous product
coordinates. Rational polyhedrality and the Euclidean Hausdorff-normalization
identification are established in the prose proof, rather than asserted as
additional Lean conclusions.

The formalization is commit \texttt{d27b977},
\href{https://github.com/Hongyu-Cheng/weak-oertel-conjecture/pull/7}{PR \#7}
in the \href{https://github.com/Hongyu-Cheng/weak-oertel-conjecture}{weak Oertel conjecture fact graph}.
The analytic result is recorded as fact \texttt{a5291bca09ba4b1e}; the formal
proof map is \texttt{lean/ONE\_OVER\_E\_AUDIT.md}.

\newpage
\head{AI disclosure}
The construction and its proof were obtained by GPT-6 Astra in Codex, using a fact
graph system similar to Danus (arXiv:2607.06447). The resulting graph, the
\href{https://github.com/Hongyu-Cheng/weak-oertel-conjecture}{weak Oertel conjecture fact graph},
contains 3503 facts, and the counterexample is a single fact in it,
\texttt{a5291bca09ba4b1e}, whose proof uses no other fact. This note was written with
AI assistance, and the authors take full responsibility for it.
\end{document}
